\usepackage[T2A]{fontenc}
\usepackage[utf8]{inputenc}
\usepackage[english,russian]{babel}
% \usepackage{cmap}
\usepackage{url}
\usepackage{booktabs}
\usepackage{nicefrac}
\usepackage{microtype}
\usepackage{lipsum}
\usepackage{graphicx}
\usepackage{subfig}
\usepackage[square,sort,comma,numbers]{natbib}
\usepackage{doi}
\usepackage{multicol}
\usepackage{multirow}
\usepackage{tabularx}

\usepackage{tikz}
\usetikzlibrary{matrix}

% Algorithms
\usepackage{algpseudocode}
\usepackage{algorithm}

%% Шрифты
\usepackage{euscript} % Шрифт Евклид
\usepackage{mathrsfs} % Красивый матшрифт
\usepackage{extsizes}

\usepackage{makecell} % diaghead in a table
\usepackage{amsmath,amsfonts,amssymb,amsthm,mathtools,dsfont}
\usepackage{icomma}

\usepackage{hyperref}
\usepackage{xcolor}

\definecolor{linkcolor}{rgb}{0.9,0,0}
\definecolor{citecolor}{rgb}{0,0.6,0}
\definecolor{urlcolor}{rgb}{0,0,1}

\hypersetup{
	unicode=true,
	pdftitle={TITLE},
	pdfauthor={AUTHOR},
	pdfkeywords={KEYWORDS},
	colorlinks=true,
	linkcolor=linkcolor,        % внутренние ссылки
	citecolor=citecolor,         % на библиографию
	urlcolor=urlcolor           % на URL
}

\graphicspath{{figs}}

\usepackage{enumitem} % Для модификаций перечневых окружений
\usepackage{etoolbox}

\makeatletter
\expandafter\patchcmd\csname\string\algorithmic\endcsname{\itemsep\z@}{\itemsep=1.5mm}{}{}
\makeatother

\usepackage{geometry}
\geometry{a4paper, top=2cm, bottom=2cm, left=2.5cm, right=1cm}
\setlength\parindent{5ex}    % Устанавливает длину красной строки 15pt
%\linespread{1.3}             % Коэффициент межстрочного интервала
\usepackage{setspace}
\usepackage[center]{titlesec} % секции посередине

%%% Убираем точки из счётчиков (чтобы \ref не добавлял точку автоматически)
\renewcommand{\thesection}{\arabic{section}}
\renewcommand{\thesubsection}{\arabic{section}.\arabic{subsection}}
\renewcommand{\thesubsubsection}{\arabic{section}.\arabic{subsection}.\arabic{subsubsection}}

%%% Возвращаем точки только в заголовки разделов (для оформления/ГОСТа)
\usepackage{titlesec} % убедитесь, что пакет загружен один раз
\titleformat{\section}{\centering\bfseries}{\thesection.}{1em}{}
\titleformat{\subsection}{\centering\bfseries}{\thesubsection.}{1em}{}
\titleformat{\subsubsection}{\centering\bfseries}{\thesubsubsection.}{1em}{}

\usepackage{icomma}
\usepackage{amsthm}
\usepackage{graphicx}
\usepackage{amssymb}
\usepackage{amsmath}
\usepackage{graphicx}
\usepackage{color}
%\usepackage{bm}
\usepackage{tabularx}
\usepackage{url}
\usepackage{multirow}
\usepackage{wrapfig}
\usepackage{caption}
\usepackage{subcaption}
\usepackage{indentfirst}
\usepackage{accents}

% Теоремы
\newtheorem{theorem}{Теорема}
\newtheorem{lemma}{Лемма}
\newtheorem{proposition}{Утверждение}
\newtheorem*{exercise}{Упражнение}
\newtheorem*{problem}{Задача}

\newtheorem{definition}{Определение}
\newtheorem*{corollary}{Следствие}
\newtheorem*{note}{Замечание}
\newtheorem*{reminder}{Напоминание}
\newtheorem*{example}{Пример}
\newtheorem*{cexample}{Контрпример}
\newtheorem*{solution}{Решение}

\renewcommand{\abstractname}{Аннотация}

% Индикаторная функция (дублирующее определение для надёжности)
\providecommand{\ind}{\mathds{1}}

% Защита от отсутствующих определений месяцев и служебных команд
% в стилях библиографии (utf8gost71u.bst использует \bblapr, \bblmar и т.п.,
% которые не всегда определяются babel-russian в полном объёме).
\AtBeginDocument{%
  \providecommand{\bbljan}{январь}%
  \providecommand{\bblfeb}{февраль}%
  \providecommand{\bblmar}{март}%
  \providecommand{\bblapr}{апрель}%
  \providecommand{\bblmay}{май}%
  \providecommand{\bbljun}{июнь}%
  \providecommand{\bbljul}{июль}%
  \providecommand{\bblaug}{август}%
  \providecommand{\bblsep}{сентябрь}%
  \providecommand{\bbloct}{октябрь}%
  \providecommand{\bblnov}{ноябрь}%
  \providecommand{\bbldec}{декабрь}%
  % Дополнительные команды, ожидаемые ГОСТ-стилем
  \providecommand{\bbland}{и}%
  \providecommand{\bbletal}{и др.}%
  \providecommand{\bblvol}{Т.}%
  \providecommand{\bblno}{№}%
  \providecommand{\bblpages}{с.}%
  \providecommand{\bblpp}{С.}%
  \providecommand{\bbleditor}{ред.}%
  \providecommand{\bbleditors}{ред.}%
  \providecommand{\bblch}{гл.}%
  \providecommand{\bbltechreport}{Технический отчёт}%
  \providecommand{\bblmastersthesis}{Магистерская диссертация}%
  \providecommand{\bblphdthesis}{Дис.\,\dots\,канд. наук}%
  \providecommand{\bblinpress}{в печати}%
  \providecommand{\bblsubmitted}{отправлено в печать}%
  \providecommand{\bblforthcoming}{принято к печати}%
  \providecommand{\bblavailablefrom}{Режим доступа}%
  \providecommand{\bblnoaddress}{Б.м.}%
  \providecommand{\bblnopublisher}{Б.и.}%
  \providecommand{\bbljour}{Журн.}%
  \providecommand{\bblOthers}{и др.}%
}